%%%%%%%%%%%%%%%%%%%%%%%%%%%%%%%%%%%%
% MA214 Lab 0 -- Welcome to the MA214 Labs
%
% Presented in the first lab meeting. Introduces the lab schedule, what
% students are required to do, and then hands off to Tutorial 0, whose
% hash they submit to Blackboard before leaving.
%%%%%%%%%%%%%%%%%%%%%%%%%%%%%%%%%%%%

%%%%%%%%%%%%%%%%%%%%%%%%%%%%%%%%%%%%
% Slide options
%%%%%%%%%%%%%%%%%%%%%%%%%%%%%%%%%%%%

% Option 1: Slides with solutions

\documentclass[slidestop,compress,mathserif]{beamer}
\newcommand{\soln}[1]{\textit{#1}}
\newcommand{\solnGr}[1]{#1}
\usepackage{booktabs}

% Option 2: Handouts without solutions

%\documentclass[11pt,containsverbatim,handout]{beamer}
%\usepackage{pgfpages}
%\pgfpagesuselayout{4 on 1}[letterpaper,landscape,border shrink=5mm]
%\newcommand{\soln}[1]{ }
%\newcommand{\solnGr}[1]{ }

%%%%%%%%%%%%%%%%%%%%%%%%%%%%%%%%%%%%
% Style
%%%%%%%%%%%%%%%%%%%%%%%%%%%%%%%%%%%%
%%%%%%%%%%%%%%%%%%%%%%%%%%%%%%%%%%%%%%%%%%%%%%%%%%%%%%%%%%%%%%%%%%%%%%%%
% MA213/214 shared slide style
%
% "Calm & Cool" palette + content-type callout boxes.
% Overhauled (2026) from the original OpenIntro/metropolis style:
%   - all OpenIntro visual branding removed (logo, grid, colors)
%   - a low-key cool palette (see colors below)
%   - tcolorbox callout boxes so students recognize content types
% See common/README.md for the command cheat-sheet.
%%%%%%%%%%%%%%%%%%%%%%%%%%%%%%%%%%%%%%%%%%%%%%%%%%%%%%%%%%%%%%%%%%%%%%%%

\usetheme{metropolis}

%%%%%%%%%%%%%%%%
% Packages
%%%%%%%%%%%%%%%%

\usepackage{geometry}
\usepackage{graphicx}
\usepackage{amssymb}
\usepackage{epstopdf}
\usepackage{amsmath}  	% this permits text in eqnarray among other benefits
\usepackage{url}		% produces hyperlinks
\usepackage{hyperref}	% allows for color usage in tables
\usepackage[english]{babel}
\usepackage[latin1]{inputenc}
\usepackage{colortbl}	% allows for color usage in tables
\usepackage{multirow}	% allows for rows that span multiple rows in tables
\usepackage{color}		% this package has a variety of color options
\usepackage{pgf}
\usepackage{calc}
\usepackage{ulem}
\usepackage{multicol}
\usepackage{textcomp}
\usepackage{txfonts}
\usepackage{listings}
\usepackage{tikz}
\usepackage{array}
\usepackage{wasysym}
\usepackage{fancyvrb}
\usepackage{etoolbox}             % \ifstrempty for optional box subtitles
\usepackage[most]{tcolorbox}      % content-type callout boxes
\tcbuselibrary{listings}          % R-code block (\begin{rblock})
\usepackage{fontawesome5}         % box icons

%%%%%%%%%%%%%%%%
% Remove navigation symbols
%%%%%%%%%%%%%%%%

\setbeamertemplate{navigation symbols}{}

%%%%%%%%%%%%%%%%
% "Calm & Cool" palette
%%%%%%%%%%%%%%%%

% chrome / structure / body
\definecolor{ccSlate}{HTML}{2E3742}   % frametitle bar, structure, section pages (deep neutral slate)
\definecolor{ccInk}{HTML}{2B2B2B}     % body text
\definecolor{ccWash}{HTML}{FCF8F1}    % gentle warm title-slide background wash (complements the cool blues)

% example (blue)
\definecolor{ccBlue}{HTML}{2F6690}
\definecolor{ccBlueBg}{HTML}{EAF1F7}
% interactive activity / discussion (green)
\definecolor{ccGreen}{HTML}{4E8A6B}
\definecolor{ccGreenBg}{HTML}{E9F3EC}
% solution (purple)
\definecolor{ccPurple}{HTML}{6F5499}
\definecolor{ccPurpleBg}{HTML}{EFEAF6}
% worksheet / focused work (terracotta)
\definecolor{ccTerra}{HTML}{C2693E}
\definecolor{ccTerraBg}{HTML}{FAEDE4}
% formula / definition (neutral slate)
\definecolor{ccGray}{HTML}{5C6B7A}
\definecolor{ccGrayBg}{HTML}{EEF1F4}
% caution / common pitfall (brick red)
\definecolor{ccRed}{HTML}{B23A48}
\definecolor{ccRedBg}{HTML}{FAEAEC}
% key idea / takeaway (muted gold)
\definecolor{ccGold}{HTML}{9C7A28}
\definecolor{ccGoldBg}{HTML}{F7F1DF}
% learning goals / lesson outcomes (teal)
\definecolor{ccTeal}{HTML}{2C6E6B}
\definecolor{ccTealBg}{HTML}{E6F0EF}
% learning objectives (indigo)
\definecolor{ccIndigo}{HTML}{45507A}
\definecolor{ccIndigoBg}{HTML}{ECEEF4}
% R code / output (gray)
\definecolor{ccCode}{HTML}{5F5E5A}
\definecolor{ccCodeBg}{HTML}{F4F5F6}

% neutral grays kept from the original style
\definecolor{gray}{rgb}{0.5, 0.5, 0.5}
\definecolor{darkGray}{rgb}{0.3, 0.3, 0.3}
\definecolor{darkerGray}{rgb}{0.2, 0.2, 0.2}

% Backwards-compatible aliases: old OpenIntro color names now point at the
% new palette, so existing \textcolor{oiB}{...}, \hl{...}, etc. keep working.
\colorlet{oiB}{ccBlue}
\colorlet{oiG}{ccGreen}
\colorlet{hlblue}{ccBlue}
\colorlet{rubineRed}{ccRed}
\colorlet{irishGreen}{ccGreen}
\colorlet{lightGreen}{ccGreen}

%%%%%%%%%%%%%%%%
% Restyle metropolis with the Calm & Cool chrome
%%%%%%%%%%%%%%%%

\metroset{block=fill}

\setbeamercolor{normal text}{fg=ccInk, bg=white}
\setbeamercolor{structure}{fg=ccSlate}
\setbeamercolor{alerted text}{fg=ccBlue}
\setbeamercolor{example text}{fg=ccGreen}
\setbeamercolor{frametitle}{bg=ccSlate, fg=white}
\setbeamercolor{progress bar}{fg=ccBlue}
\setbeamercolor{title separator}{fg=ccBlue}
\setbeamercolor{title}{fg=ccSlate}
\setbeamercolor{subtitle}{fg=ccBlue}
\setbeamercolor{author}{fg=ccInk}
\setbeamercolor{institute}{fg=ccGray}
\setbeamercolor{date}{fg=ccGray}
\setbeamercolor{section title}{fg=ccSlate}
\setbeamercolor{standout}{fg=white, bg=ccSlate}

% Section pages sit on the same warm wash as the title slide. We keep
% metropolis's section-title styling and just paint a full-page ccWash
% background behind it (needs >=2 pdflatex passes for the current-page node).
\setbeamertemplate{section page}{%
  \begin{tikzpicture}[remember picture, overlay]
    \fill[ccWash] (current page.south west) rectangle (current page.north east);
  \end{tikzpicture}%
  \centering
  \begin{minipage}{22em}
    \raggedright
    \usebeamercolor[fg]{section title}%
    \usebeamerfont{section title}%
    \insertsectionhead\\[-1ex]
    \usebeamertemplate*{progress bar in section page}%
    \par
  \end{minipage}\par
  \vspace{\baselineskip}
}

%%%%%%%%%%%%%%%%
% Custom inline text commands (recolored to the palette)
%%%%%%%%%%%%%%%%

% degree
\newcommand{\degree}{\ensuremath{^\circ}}

% cite
\newcommand{\ct}[1]{
\vfill
{\tiny #1}}

\newcommand{\NamedNote}[2]{
\rule{2.5cm}{0.25pt} \\ \textit{\footnotesize{\textcolor{rubineRed}{#1:} \textcolor{darkerGray}{#2}}}}

% Note
\newcommand{\Note}[1]{
\rule{2.5cm}{0.25pt} \\ \textit{\footnotesize{\textcolor{rubineRed}{Note:} \textcolor{darkerGray}{#1}}}}

% Remember
\newcommand{\Remember}[1]{\textit{\scriptsize{\textcolor{ccTerra}{Remember:} #1}}}

% expected counts
\newcommand{\ex}[1]{\textit{\textcolor{ccBlue}{#1}}}

% red
\newcommand{\red}[1]{\textit{\textcolor{rubineRed}{#1}}}

% pink
\newcommand{\pink}[1]{\textit{\textcolor{rubineRed!90!white!50}{#1}}}

% green
\newcommand{\green}[1]{\textit{\textcolor{irishGreen}{#1}}}

% orange
\newcommand{\orange}[1]{\textit{\textcolor{ccTerra}{#1}}}

% links: webURL, webLink, appLink
\newcommand{\webURL}[1]{\urlstyle{same}{ \textit{\textcolor{darkGray}{\url{#1}}}}}
\newcommand{\webLink}[2]{\href{#1}{\textcolor{darkGray}{{#2}}}}
\newcommand{\appLink}[2]{\href{#1}{\textcolor{white}{{#2}}}}

% mail
\newcommand{\mail}[1]{\href{mailto:#1}{\textit{\textcolor{darkGray}{#1}}}}

% highlighting: hl, hlGr, mathhl
\newcommand{\hl}[1]{\textit{\textcolor{hlblue}{#1}}}
\newcommand{\hlGr}[1]{\textit{\textcolor{lightGreen}{#1}}}
\newcommand{\mathhl}[1]{\textcolor{hlblue}{\ensuremath{#1}}}

%%%%%%%%%%%%%%%%
% Structural helpers (unchanged from the original)
%%%%%%%%%%%%%%%%

% two col: two columns
\newenvironment{twocol}[4]{
\begin{columns}[c]
\column{#1\textwidth}
#3
\column{#2\textwidth}
#4
\end{columns}
}

% slot (for probability calculations)
\newenvironment{slot}[2]{
\begin{array}{c}
\underline{#1} \\
#2
\end{array}
}

% pr: left and right parentheses
\newcommand{\pr}[1]{
\left( #1 \right)
}

% cancel
\newcommand{\cancel}[1]{%
    \tikz[baseline=(tocancel.base)]{
        \node[inner sep=0pt,outer sep=0pt] (tocancel) {#1};
        \draw[red, line width=0.5mm] (tocancel.south west) -- (tocancel.north east);
    }%
}

% removepagenumbers
\newcommand{\removepagenumbers}{%
  \setbeamertemplate{footline}{}
}

% change margin
\newenvironment{changemargin}[2]{%
\begin{list}{}{%
\setlength{\topsep}{0pt}%
\setlength{\leftmargin}{#1}%
\setlength{\rightmargin}{#2}%
\setlength{\listparindent}{\parindent}%
\setlength{\itemindent}{\parindent}%
\setlength{\parsep}{\parskip}%
}%
\item}{\end{list}}

% footnote
\long\def\symbolfootnote[#1]#2{\begingroup%
\def\thefootnote{\fnsymbol{footnote}}\footnote[#1]{#2}\endgroup}

% commands from the book
\newenvironment{data}[1]{\texttt{#1}}{}
\newenvironment{var}[1]{\texttt{#1}}{}
\newenvironment{resp}[1]{\texttt{#1}}{}

%%%%%%%%%%%%%%%%
% Content-type callout boxes (tcolorbox)
%
% Each type has a colored title strip (icon + label) over a light fill.
% Color = content type; the label + icon reinforce it (so the cue does
% not rely on color alone).
%%%%%%%%%%%%%%%%

\tcbset{
  cccallout/.style 2 args={
    enhanced, breakable=false,
    colback=#2, colframe=#1, colbacktitle=#1, coltitle=white,
    fonttitle=\bfseries\footnotesize,
    boxrule=0.6pt, arc=2.5pt, outer arc=2.5pt,
    left=5pt, right=5pt, top=3pt, bottom=3pt, boxsep=2pt,
    toptitle=2pt, bottomtitle=2pt, lefttitle=5pt,
    before skip=5pt, after skip=5pt,
  },
}

% One-off / collection of examples (blue): \workedexample[subtitle]{body}.
% (Named \workedexample, not \example, because beamer defines \example.)
\newtcolorbox{examplebox}[1][]{cccallout={ccBlue}{ccBlueBg},
  title={\faIcon{lightbulb}~~Example\ifstrempty{#1}{}{ \textmd{---}\ #1}}}
\newcommand{\workedexample}[2][]{\begin{examplebox}[#1]#2\end{examplebox}}

% Running-example header (blue): \examplehead{name} at the top of each frame of a
% multi-slide running example. A one-off or a collection of examples uses the
% \workedexample box above instead.
\newcommand{\examplehead}[1]{%
  \vspace{-0.4em}\par\noindent\tcbox[on line, colback=ccBlue, colframe=ccBlue, coltext=white,
    boxrule=0pt, arc=2pt, left=5pt, right=5pt, top=1pt, bottom=1pt,
    box align=base, nobeforeafter]{\footnotesize\bfseries\faIcon{lightbulb}\,Running example}%
  \ifstrempty{#1}{}{\enspace{\footnotesize\bfseries\textcolor{ccBlue}{#1}}}%
  \par\nobreak\vspace{-0.6em}{\color{ccBlue}\rule{\linewidth}{0.7pt}}\par\nobreak\vspace{-0.8em}}

% Interactive activity (green): \activity[subtitle]{body}
\newtcolorbox{activitybox}[1][]{cccallout={ccGreen}{ccGreenBg},
  title={\faIcon{users}~~Activity\ifstrempty{#1}{}{ #1}}}
\newcommand{\activity}[2][]{\begin{activitybox}[#1]#2\end{activitybox}}

% Discussion question (green): \dq{body}
\newtcolorbox{dqbox}{cccallout={ccGreen}{ccGreenBg}, title={\faIcon{comments}~~Discussion}}
\newcommand{\dq}[1]{\begin{dqbox}#1\end{dqbox}}

% Solution (purple): \solnbox{body}. Gated by the per-deck \soln toggle, so it
% shows in the instructor build and vanishes from the student handout build.
% (Named \solnbox, not \solution, because beamer already defines \solution.)
\newtcolorbox{solutionbox}{cccallout={ccPurple}{ccPurpleBg}, title={\faIcon{key}~~Solution}}
\newcommand{\solnbox}[1]{\soln{\begin{solutionbox}#1\end{solutionbox}}}

% Worksheet / focused work (terracotta): \worksheet[subtitle]{body}
\newtcolorbox{worksheetbox}[1][]{cccallout={ccTerra}{ccTerraBg},
  title={\faIcon{pencil-alt}~~Worksheet\ifstrempty{#1}{}{ #1}}}
\newcommand{\worksheet}[2][]{\begin{worksheetbox}[#1]#2\end{worksheetbox}}

% Formula / definition (neutral slate). Supports both real usages found in the
% decks: \formula{title}{body} and the bare \formula{body} (title defaults to
% "Formula"). The second group is an optional brace argument.
\newtcolorbox{formulabox}[1][Formula]{cccallout={ccGray}{ccGrayBg},
  title={\faIcon{square-root-alt}~~#1}}
\NewDocumentCommand{\formula}{m g}{%
  \IfNoValueTF{#2}%
    {\begin{formulabox}#1\end{formulabox}}%
    {\begin{formulabox}[#1]#2\end{formulabox}}}

% Common pitfall / caution (brick red): \caution[subtitle]{body}
\newtcolorbox{cautionbox}[1][]{cccallout={ccRed}{ccRedBg},
  title={\faIcon{exclamation-triangle}~~Caution\ifstrempty{#1}{}{ #1}}}
\newcommand{\caution}[2][]{\begin{cautionbox}[#1]#2\end{cautionbox}}

% Key idea / takeaway (muted gold): \keyidea[subtitle]{body}
\newtcolorbox{keyideabox}[1][]{cccallout={ccGold}{ccGoldBg},
  title={\faIcon{star}~~Key idea\ifstrempty{#1}{}{ #1}}}
\newcommand{\keyidea}[2][]{\begin{keyideabox}[#1]#2\end{keyideabox}}

% Learning goals / lesson outcomes (teal): \goals{body}. Sits under the short
% LO headings on the objectives slide; body is the "you should be able to" list.
\newtcolorbox{goalsbox}{cccallout={ccTeal}{ccTealBg}, fontupper=\footnotesize,
  top=1pt, bottom=2pt, boxsep=1pt, before skip=3pt,
  title={\faIcon{bullseye}~~By the end of today, you should be able to\ldots}}
\newcommand{\goals}[1]{\begin{goalsbox}#1\end{goalsbox}}

% Learning objectives (indigo): \objectives{body} -- the formal LO headings,
% shown alongside \goals on the objectives slide.
\newtcolorbox{objectivesbox}{cccallout={ccIndigo}{ccIndigoBg}, fontupper=\small,
  top=1pt, bottom=2pt, boxsep=1pt, before skip=3pt,
  title={\faIcon{clipboard-list}~~Learning objectives}}
\newcommand{\objectives}[1]{\begin{objectivesbox}#1\end{objectivesbox}}

% Defaults for a bare \begin{lstlisting}. Without these it inherits the listings
% package defaults: 11pt type, fixed-width columns (which spaces code out as
% "a g e _ f i t"), and no line breaking, so a long line silently runs off the
% right edge of the slide. A deck that passes its own basicstyle still wins.
% Layout only -- syntax coloring is left to the \begin{rblock} environment.
\lstset{
  basicstyle=\ttfamily\footnotesize, columns=fullflexible,
  breaklines=true, breakindent=1.5em, showstringspaces=false,
  numbers=none, frame=none}

% R code / output (gray): \begin{rblock}[subtitle] ... \end{rblock}
% NOTE: a frame containing an rblock must be declared \begin{frame}[fragile].
\lstdefinestyle{ccR}{
  language=R, basicstyle=\ttfamily\footnotesize, columns=fullflexible,
  breaklines=true, showstringspaces=false, numbers=none, frame=none,
  keywordstyle=\color{ccBlue}, commentstyle=\color{ccGreen!75!black},
  stringstyle=\color{ccTerra}}
\newtcblisting{rblock}[1][]{cccallout={ccCode}{ccCodeBg},
  title={\faIcon{terminal}~~R\ifstrempty{#1}{}{ #1}},
  listing only, listing style=ccR}

% --- Legacy boxes (kept so existing decks compile) ---

% app: application exercise -> alias of activity
\newcommand{\app}[1]{\activity{#1}}

% pq: practice question (green, interactive family)
\newtcolorbox{pqbox}{cccallout={ccGreen}{ccGreenBg}, title={\faIcon{list-ol}~~Practice}}
\newcommand{\pq}[1]{\begin{pqbox}#1\end{pqbox}}

% solnMult: solutions for multiple-choice practice questions (uses \soln toggle)
\newcommand{\solnMult}[1]{
\item[] \vspace{-0.59cm}
\only<1>{\item #1}
\soln{\only<2->{\item \orange{#1}}}
}

%%%%%%%%%%%%%%%%
% De-branded title frame + attribution
%%%%%%%%%%%%%%%%

% Eyebrow line on the title slide; a deck may \renewcommand this.
\newcommand{\coursetag}{MA214: Applied Statistics}

% Attribution: all slides released under CC BY-SA, crediting both the
% instructor (later material) and OpenIntro (basis of the earlier material).
\newcommand{\courseattribution}{%
  Slides by Jonathan H.~Huggins, building on materials from OpenIntro's
  \emph{Introduction to Modern Statistics} (Mine \c{C}etinkaya-Rundel et al.).
  Released under the
  \webLink{http://creativecommons.org/licenses/by-sa/3.0/us/}{CC BY-SA license}.
  Some images included under fair use (educational purposes).}

% \maketitleframe: de-branded title slide (no logo, no grid). Uses the deck's
% \title, \subtitle, and \coursetag.
\newcommand{\maketitleframe}{%
  \addtocounter{framenumber}{-1}%
  {\removepagenumbers
  \setbeamercolor{background canvas}{bg=ccWash}%
  \begin{frame}[plain]
    \vspace*{2.6em}
    {\color{ccSlate}\rule{\textwidth}{2.5pt}}\par\vspace{1.5em}
    {\usebeamerfont{date}\scshape\color{ccGray}\coursetag}\par\vspace{0.55em}
    {\usebeamerfont{title}\color{ccSlate}\inserttitle\par}
    \vspace{0.45em}{\color{ccBlue}\rule{3em}{2.5pt}}\par\vspace{0.7em}
    {\usebeamerfont{subtitle}\color{ccBlue}\insertsubtitle\par}
    \vfill
    {\color{ccGray!55}\rule{\textwidth}{0.4pt}}\par\vspace{0.45em}
    {\tiny\color{ccGray}\courseattribution\par}
    \vspace*{0.4em}
  \end{frame}}}

%%%%%%%%%%%%%%%%
% Graphics
%%%%%%%%%%%%%%%%

\DeclareGraphicsRule{.tif}{png}{.png}{`convert #1 `dirname #1`/`basename #1 .tif`.png}

\graphicspath{ {../../common/} }

% This is a lab deck, not a lecture deck: credit the lab instructor.
\renewcommand{\coursetag}{MA214: Applied Statistics --- Lab}
\renewcommand{\courseattribution}{%
  MA214 lab materials by Yongho Lim. Released under the
  \webLink{http://creativecommons.org/licenses/by-sa/3.0/us/}{CC BY-SA license}.}

% EDIT ME each semester: the Docker Hub account hosting the tutorial image
% is written literally in the "Getting the tutorials running" frame below
% (the code block is verbatim, so a macro would not expand there).

%%%%%%%%%%%%%%%%%%%%%%%%%%%%%%%%%%%%%%%%%%%%%%%%%%%%%%%%%%%%%%%%%%%%%%%%
% Preamble
%%%%%%%%%%%%%%%%%%%%%%%%%%%%%%%%%%%%%%%%%%%%%%%%%%%%%%%%%%%%%%%%%%%%%%%%

\title[Lab 0]{Lab 0: Welcome to the labs}
\subtitle{Module 1: Modeling and Linear Regression}
\date{}

%%%%%%%%%%%%%%%%%%%%%%%%%%%%%%%%%%%%%%%%%%%%%%%%%%%%%%%%%%%%%%%%%%%%%%%%
% Begin document
%%%%%%%%%%%%%%%%%%%%%%%%%%%%%%%%%%%%%%%%%%%%%%%%%%%%%%%%%%%%%%%%%%%%%%%%

\begin{document}


%%%%%%%%%%%%%%%%%%%%%%%%%%%%%%%%%%%%%%%%%%%%%%%%%%%%%%%%%%%%%%%%%%%%%%%%
% Title page
%%%%%%%%%%%%%%%%%%%%%%%%%%%%%%%%%%%%%%%%%%%%%%%%%%%%%%%%%%%%%%%%%%%%%%%%

\maketitleframe

% !TEX root =  Lab0.tex


%%%%%%%%%%%%%%%%%%%%%%%%%%%%%%%%%%%%
% Recap/Agenda
%%%%%%%%%%%%%%%%%%%%%%%%%%%%%%%%%%%%
\begin{frame}
\frametitle{MA214 Labs}
    \begin{itemize}
    \item \hl{Previously:} nothing --- this is the first lab meeting.
    \item \hl{This time:} the machinery of the labs, so that from Lab 1
      onward none of it is new.
      \begin{itemize}
      \item the two kinds of lab, and the semester schedule
      \item what a lab must contain to count as \hl{completed}
      \item the tutorial system, end to end: run it, work it, submit it
      \end{itemize}
    \item \hl{Reading:} the syllabus, lab section.
    \item \hl{Due today:} the \hl{Tutorial 0} hash, on Blackboard before you leave.
    \item \hl{Due before Lab 1:} the Lab 1 tutorial.
    \end{itemize}

\vspace{0.4em}
No statistics today.
\end{frame}


%%%%%%%%%%%%%%%%%%%%%%%%%%%%%%%%%%%%
% Learning objectives:
%%%%%%%%%%%%%%%%%%%%%%%%%%%%%%%%%%%%
\begin{frame}
    \frametitle{Lab Objectives}
    \goals{%
    \begin{itemize}
    \item say what happens in a lab, and which labs are which;
    \item state exactly what you must submit for a lab to count as complete;
    \item open a tutorial and run code inside it; and
    \item generate a submission hash and post it to Blackboard.
    \end{itemize}}
    \objectives{%
    \begin{itemize}
    \item \textbf{(Lab policy)} What a completed lab requires --- attendance,
      tutorial hash, worksheet
    \item \textbf{(Lab tooling)} Running the tutorial container and submitting work
    \end{itemize}}
\end{frame}


%%%%%%%%%%%%%%%%%%%%%%%%%%%%%%%%%%%%%%%%%%%%%%%%%%%%%%%%%%%%%%%%%%%%%%%%
\section{How the labs run}
%%%%%%%%%%%%%%%%%%%%%%%%%%%%%%%%%%%%%%%%%%%%%%%%%%%%%%%%%%%%%%%%%%%%%%%%
\begin{frame}
\frametitle{Who is in the room}

\begin{itemize}
  \item \hl{Lab instructor:} Dr.~Yongho Lim --- \mail{ylim2@bu.edu} \\
        Office hours: Monday 1:00--2:00pm, Thursday 12:00--1:00pm, CDS 408
  \item \hl{Teaching fellows:} Sidharth Soundararajan (\mail{sidsound@bu.edu})
  \item \hl{Learning assistants:} Cynthia Huang (\mail{chuang12@bu.edu})

\end{itemize}

\vspace{0.5em}
Teaching fellow and learning assistant office hours: TBA, posted on the
course page.


\end{frame}
%%%%%%%%%%%%%%%%%%%%%%%%%%%%%%%%%%%%%%%%%%%%%%%%%%%%%%%%%%%%%%%%%%%%%%%%
\begin{frame}
\frametitle{What the labs are for}

Lecture is where the statistics is developed. Lab is where you
\hl{do} it: on real data, in R, with other people.

\vspace{0.5em}
\begin{itemize}
  \item You work in \hl{groups}. Statistics is argued about, not
        recited.
  \item You use \hl{R} and RStudio. The code is a means, not the point.
  \item You produce something every session --- a worksheet, a
        deliverable.
\end{itemize}

\keyidea{The lab grade is not about getting clean output. It is about
whether you can interpret what the output says and say why it is
reasonable.}

\end{frame}
%%%%%%%%%%%%%%%%%%%%%%%%%%%%%%%%%%%%%%%%%%%%%%%%%%%%%%%%%%%%%%%%%%%%%%%%
\begin{frame}
\frametitle{Two kinds of lab}

\begin{columns}[T]
\column{0.48\textwidth}
\textbf{\color{ccBlue} Skills labs (5)}

\vspace{0.3em}
Short, self-contained, hands-on.

\begin{itemize}\footnotesize
\item Prepare with an online tutorial
\item Work an in-lab worksheet
\item Graded \emph{completed} / \emph{not completed}
\end{itemize}

\column{0.48\textwidth}
\textbf{\color{ccTerra} Project labs (2 projects)}

\vspace{0.3em}
Several weeks each, one group product.

\begin{itemize}\footnotesize
\item Project 1: video presentation
\item Project 2: written report
\item Graded E / S / R / N
\end{itemize}
\end{columns}

\vspace{0.8em}
Roughly half the semester is skills labs, half is project work. Both
count toward your final grade, and they count in different ways --- see
the syllabus grading table.

\end{frame}
%%%%%%%%%%%%%%%%%%%%%%%%%%%%%%%%%%%%%%%%%%%%%%%%%%%%%%%%%%%%%%%%%%%%%%%%
\begin{frame}
\frametitle{The lab schedule}

\vspace{-0.8em}
{\scriptsize
\renewcommand{\arraystretch}{0.9}
\begin{center}
\begin{tabular}{cll}
\toprule
\bf Week & \bf Session & \bf Focus \\
\midrule
1 & \color{ccBlue}Lab 0  & Welcome to the labs \\
2  & \color{ccBlue}Lab 1  & Working with R and data \\
3  & \color{ccBlue}Lab 2  & Regression models \\
4  & \color{ccTerra}P1--1 & Project 1 launch and outline \\
5  & \color{ccTerra}P1--2 & Project 1 data cleaning and EDA \\
6  & \color{ccTerra}P1--3 & Project 1 modeling and diagnostics \\
7  & \color{ccTerra}P1--4 & Project 1 presentation work session \\
8  & \color{ccTerra}P2--1 & Project 2 launch: bootstrapping \\
9  & \color{ccTerra}P2--2 & Project 2 simulation design \\
10  & \color{ccTerra}P2--3 & Project 2 real data analysis \\
11 & \color{ccTerra}P2--4 & Project 2 final work session \\
12 & \color{ccBlue}Lab 3  & Bayes' rule \\
13 & \color{ccBlue}Lab 4  & Conjugate families: beta-binomial \\
14 & \color{ccGray}---    & \color{ccGray}No lab \\
15 & \color{ccBlue}Lab 5  & Bayesian linear regression \\
\bottomrule
\end{tabular}
\end{center}}

\vspace{-0.6em}
\Note{Calendar dates follow your own section's meeting day --- check
your schedule and the course page.}

\end{frame}
%%%%%%%%%%%%%%%%%%%%%%%%%%%%%%%%%%%%%%%%%%%%%%%%%%%%%%%%%%%%%%%%%%%%%%%%
\section{What counts as a completed lab}
%%%%%%%%%%%%%%%%%%%%%%%%%%%%%%%%%%%%%%%%%%%%%%%%%%%%%%%%%%%%%%%%%%%%%%%%
\begin{frame}
\frametitle{The shape of a skills lab}

Every one of the five runs the same way.

\begin{enumerate}
  \item \hl{Before lab.} Work the online \hl{tutorial} for that lab.
        It practises the R on a different scenario, so that syntax is not
        what slows you down in the room.
  \item \hl{In lab.} Work the \hl{worksheet} with your group, with the
        lab team circulating.
  \item \hl{Submit.} The tutorial's \hl{hash} to Blackboard, and the
        worksheet as instructed for that lab.
\end{enumerate}

\keyidea{A skills lab counts as completed only when \emph{both} arrive:
the tutorial hash \hl{and} the worksheet. One without the other is not.}

\end{frame}
%%%%%%%%%%%%%%%%%%%%%%%%%%%%%%%%%%%%%%%%%%%%%%%%%%%%%%%%%%%%%%%%%%%%%%%%
\begin{frame}
\frametitle{What is required of you}

\begin{itemize}
  \item \hl{Show up.} Lab attendance is graded. For an A you may miss
        at most one lab.
  \item \hl{Come prepared.} Do the tutorial \emph{before} the lab it
        belongs to, not after.
  \item \hl{Submit both pieces} for each skills lab: hash and worksheet.
  \item \hl{Carry your weight on projects.} Every deliverable is
        followed by a peer assessment form.
\end{itemize}

\caution[--- absences]{An absence is excused only if you email the lab
instructor \emph{in advance} with your reason. Excused absences are not
counted as missed. After the fact is too late.}

\end{frame}
%%%%%%%%%%%%%%%%%%%%%%%%%%%%%%%%%%%%%%%%%%%%%%%%%%%%%%%%%%%%%%%%%%%%%%%%
\begin{frame}
\frametitle{The lab projects, in one slide}

You will work in the same group on two projects.

\vspace{0.5em}
\begin{itemize}
  \item \hl{Project 1} --- a video presentation. Weeks 4--7.
  \item \hl{Project 2} --- a written report. Weeks 8--11.
\end{itemize}

\vspace{0.5em}
Each is graded on the ESRN rubric: \hl{E}xcellent, \hl{S}atisfactory,
\hl{R}evisions needed, \hl{N}ot assessable. An S earns full credit. An R
or N must be revised and resubmitted \emph{once}.

\vspace{0.5em}
Both projects must reach S or E by the end of the course.

\Note{More detail in the Week 4 lab, when Project 1 launches. Nothing to
do about it today.}

\end{frame}
%%%%%%%%%%%%%%%%%%%%%%%%%%%%%%%%%%%%%%%%%%%%%%%%%%%%%%%%%%%%%%%%%%%%%%%%
\section{The tutorial system}
%%%%%%%%%%%%%%%%%%%%%%%%%%%%%%%%%%%%%%%%%%%%%%%%%%%%%%%%%%%%%%%%%%%%%%%%
\begin{frame}[fragile]
\frametitle{Getting the tutorials running}

The tutorials run in \hl{Docker}, so you install no R packages and set no
working directory. Install Docker Desktop once, then for every tutorial:

Run the tutorial launcher that is available on the blackboard. 

Then open \webURL{http://localhost:3838} in a browser (your launcher will open it for you). 
You will see a list of tutorials. Click the one you want to run.

\vspace{0.4em}
\begin{itemize}\footnotesize
  \item \hl{Back to Main Tutorial Page} at the top of a tutorial returns
        to the picker.
  \item Nothing is downloaded and nothing is installed into R --- the
        data lives inside the tutorial.
\end{itemize}

\end{frame}
%%%%%%%%%%%%%%%%%%%%%%%%%%%%%%%%%%%%%%%%%%%%%%%%%%%%%%%%%%%%%%%%%%%%%%%%
\begin{frame}
\frametitle{Inside a tutorial}

\begin{itemize}
  \item \hl{Exercise box} --- an editable code editor with a
        \emph{Run Code} button. Run it as often as you like; you cannot
        break anything.
  \item \hl{Blanks} --- \texttt{\_\_\_} marks something for you to
        replace.
  \item \hl{Hint} --- available from the start. Use it.
  \item \hl{Solution} --- released on a timer, normally \emph{after}
        that tutorial's deadline. Do not wait for it.
  \item \hl{Quiz} --- multiple choice, with retries.
\end{itemize}

\caution[--- one box, one world]{Each exercise box starts fresh. An
object you create in one box does not exist in the next one.}

\end{frame}
%%%%%%%%%%%%%%%%%%%%%%%%%%%%%%%%%%%%%%%%%%%%%%%%%%%%%%%%%%%%%%%%%%%%%%%%
\begin{frame}
\frametitle{Submitting a hash}

At the bottom of every tutorial is a \hl{Submit your work} section.

\vspace{0.6em}
\begin{enumerate}
  \item Check that you have attempted every exercise and quiz.
  \item Type your \hl{name} and \hl{student ID}.
  \item Click \hl{Generate}.
  \item Select the \hl{whole} hash and copy it.
  \item Paste it into that tutorial's submission form on
        \hl{Blackboard}.
\end{enumerate}

\vspace{0.4em}
\begin{itemize}\footnotesize
  \item No hash appears? Answer at least one question first, and check
        that neither name box is blank.
  \item Changed an answer afterwards? Generate a new hash and submit the
        newer one.
  \item The hash encodes your answers and your ID. Do not post it
        anywhere public.
\end{itemize}

\end{frame}



%%%%%%%%%%%%%%%%%%%%%%%%%%%%%%%%%%%%%%%%%%%%%%%%%%%%%%%%%%%%%%%%%%%%%%%%
\begin{frame}
\frametitle{When you are stuck}

In this order:

\begin{enumerate}
  \item Read the error message, then read the question again.
  \item Check spelling, capitalisation, brackets, quotation marks.
  \item Compare your code with the worked example above it.
  \item Open the hint.
  \item Ask your group, then an LA, a TF, or the lab instructor.
\end{enumerate}

Say what you already tried, and do not delete the broken attempt --- it
is the most useful thing you can show us.

\keyidea{An error message is not a verdict on you. It is R telling you
what it needs next.}

\end{frame}
%%%%%%%%%%%%%%%%%%%%%%%%%%%%%%%%%%%%%%%%%%%%%%%%%%%%%%%%%%%%%%%%%%%%%%%%
\section{Wrap-up}
%%%%%%%%%%%%%%%%%%%%%%%%%%%%%%%%%%%%%%%%%%%%%%%%%%%%%%%%%%%%%%%%%%%%%%%%
\begin{frame}
\frametitle{Key takeaways}

\begin{itemize}
\item Five \hl{skills labs} and two \hl{projects}, roughly half the semester each.
\item A skills lab is completed only when \emph{both} the tutorial hash and the
  worksheet arrive --- one without the other is not.
\item The tutorial for a lab is done \emph{before} that lab, not after.
\item Attendance is graded, and an absence is excused only if you email
  \hl{in advance}.
\item Tutorials run in \hl{Docker}. Download \hl{Docker Desktop} and open it. 
\item \hl{Today:} finish Tutorial 0 and submit its hash to Blackboard before
  you leave.
\item Go over the R onboarding website for R and RStudio if you are new to R.
\item https://bu-intro-stats.github.io/MA213-214-R-onboarding-website/
\end{itemize}

\end{frame}
%%%%%%%%%%%%%%%%%%%%%%%%%%%%%%%%%%%%%%%%%%%%%%%%%%%%%%%%%%%%%%%%%%%%%%%%

%%%%%%%%%%%%%%%%%%%%%%%%%%%%%%%%%%%%
% End document
%%%%%%%%%%%%%%%%%%%%%%%%%%%%%%%%%%%%

\end{document}
